\subsection{Observación, presentación y recepción de pilotos}
La ejecución directa del registro se distingue del período observado. IN1 admite solicitudes ingresadas durante los primeros 21 días naturales del mes y las sigue durante siete días desde su envío; el último seguimiento termina el día 28. Se conservan las solicitudes sin respuesta y las condiciones de comparabilidad. IN4 mide el tiempo activo de las solicitudes de ese mismo universo, con indicadores y exclusiones propios. IN5 conserva 28 días naturales de base y otros 28 de evaluación, con la acción y sus evidencias entre ambos. Los paquetes quincenales realizan captura y control durante la ventana; no acreditan su cierre definitivo al consumir sus HH antes del último dato.

El registro \path{ventanas_observacion.csv} sitúa las bases en 1--28 de junio de 2028 y las evaluaciones en 1--28 de julio de 2028, para el escenario de inicio de diciembre de 2026 y la programación de origen de SD13. Las acciones IN5 se ubican en 29--30 de junio, con acceso permitido y habilitación validada. La formación IN1 depende del cierre de su base y debe terminar antes de activar la cohorte del piloto. Sus sesiones y evaluación requieren una asignación específica en los días restantes; las HH del paquete no prueban por sí solas que esa ventana sea suficiente.

Calidad emitirá los informes definitivos después del cierre de los datos, conservando versiones, participantes, exclusiones, casos pendientes e indicadores. La contraparte dispondrá de diez días hábiles de revisión y, si formula observaciones, Synaptix dispondrá de diez días hábiles para subsanarlas. La subsanación requiere carga y ejecutores identificados; una ventana temporal no aporta HH de corrección implícitas. La repetición de observaciones de la misma naturaleza mantiene el tratamiento de atraso previsto en las bases \parencite[art.~18.2--18.4; p.~13]{bases-admin}.

\begingroup\small\setlength{\tabcolsep}{3pt}
\begin{longtable}{L{41mm}L{24mm}L{24mm}L{24mm}L{24mm}}
\caption{Recepción de pilotos y margen hasta H12 en el escenario de inicio diciembre de 2026}\label{tab:sd7-recepcion-pilotos}\\\hline
Escenario & Informe & Fin revisión & Fin subsanación & Transferencia\\\hline\endfirsthead
\hline Escenario & Informe & Fin revisión & Fin subsanación & Transferencia\\\hline\endhead
Reparto directo secuencial & 03/08/2028 & 18/08/2028 & 01/09/2028 & 07/09/2028\\\hline
Preparación previa y cierre paralelo de informes & 30/07/2028 & 11/08/2028 & 28/08/2028 & 30/08/2028\\\hline
Anticipación de base y piloto a meses 18/19 & 04/07/2028 & 18/07/2028 & 01/08/2028 & 03/08/2028\\\hline
\end{longtable}\FuenteTabla{Escenarios de Synaptix: cierre de observación 28/07/2028 en las dos primeras filas y 28/06/2028 en el anticipo; revisión y subsanación diez días hábiles cada una; lunes a viernes, con 15 de agosto bloqueado por hipótesis conservadora. H12 se contrasta con 31/08/2028; se requiere validar feriados y calendario efectivo.}\endgroup


La subred \path{red_observacion_recepcion.csv} conserva los eventos de cierre, presentación, revisión, subsanación y validación final de transferencia. El primer escenario adopta cuatro días hábiles para terminar un informe de 32 HH conforme a la duración directa de T-15, después del cierre de observación. Con la reserva completa de subsanación y transferencia secuencial, termina el 7 de septiembre, fuera de H12. La programación no puede considerar el informe concluido en una quincena anterior a su último dato.

El segundo escenario estudia preparación previa de insumos y dos frentes de cierre del informe durante el 29--30 de julio. Conserva sus 32 HH y reparte la contribución de Calidad de 16 HH en dos revisiones de 8 HH, sobre secciones identificadas y coordinadas por el responsable de Calidad. El trabajo de fin de semana requiere turnos, descansos y disponibilidad reservados; no se atribuye a una autorización o ejecución realizadas. La documentación de transferencia se prepara durante la revisión y la validación final posterior conserva ocho HH de Calidad en dos días hábiles. Una corrección que invalide esa preparación obliga a recalcular su esfuerzo y duración. Este escenario no sustituye la asignación activa de T-15 hasta incorporar esos ejecutores e intervalos.

Incluso ese cierre paralelo deja sólo el 31 de agosto antes de H12 después de agotar revisión y subsanación. No absorbe los diez días hábiles de reserva marina planteados por campaña. La corrección documental y las cancelaciones por marejada no se imputan a una misma reserva como si nunca pudieran coincidir. El cálculo conserva lunes--viernes y bloquea el 15 de agosto por hipótesis conservadora; se completará con feriados y ventanas efectivas. Ninguno de estos dos escenarios acredita suficiencia integral.

\paragraph{Objetivo de anticipación y control de la línea base.}
El recorrido hacia atrás desde el 31 de agosto de 2028, reservando diez días hábiles de revisión, diez de subsanación, dos de validación final y diez de reserva marina, exige presentar el expediente a más tardar el 17 de julio en este escenario. El cierre de observación del 28 de julio impide cumplir ese objetivo con la programación de origen, aunque se agreguen revisores. Proyecto reprogramará la cadena completa de preparación, construcción, pruebas, base y piloto de las innovaciones afectadas, conservando sus duraciones y condiciones de evaluación, y reconciliará las fechas con SD13, T-14, T-15 y T-18 antes de aprobarla. Anticipar sólo la fecha del informe sin adelantar sus datos no es una solución admisible.

La anticipación se desarrolla trasladando base y evaluación a los meses 18 y 19 y adelantando un mes los paquetes de innovación de 14--20, mientras la materialización contractual conserva el mes 21. El ensayo de programación conserva las 20.820 HH, los 425 enlaces resueltos entre paquetes fijos y el pico de 35 refuerzos en el mes 8. No acredita las puertas simbólicas, las habilitaciones ni el calendario efectivo. Se conserva como alternativa reproducible en \path{ventanas_anticipo_propuesto.csv}, pendiente de incorporar los eventos reales de formación y cierre y reconciliar el calendario de SD13.

En esta alternativa las bases cierran el 28 de mayo, la evaluación el 28 de junio y el informe se presenta el 4 de julio bajo las hipótesis del cálculo. Revisión y subsanación concluyen el 1 de agosto; la validación final termina el 3 de agosto. Los diez días hábiles de reserva marina terminan el 18 de agosto, antes de H12. Este margen temporal permite continuar el diseño sin acortar la observación ni las revisiones. Requiere verificar que la construcción y las pruebas anticipadas estén liberadas y que la capacitación posterior a la base quepa en los turnos previos al 1 de junio; no se atribuye aceptación por consumir ese margen. Hasta completar esa comprobación, la asignación activa conserva sus resultados de esfuerzo y queda condicionada por estos eventos temporales.

IN3 fijará antes de medir la ventana de muestra manual/asistida, su último caso y el cierre de revisión humana; su informe dependerá de esos eventos y de su recepción, sin atribuirle automáticamente las ventanas de IN1 o IN5. IN2 mantiene sus cierres estacionales, revisiones por lotes antes de botaduras y aceptación propia. Los informes de innovación no sustituyen la verificación del volumen real durante las cuatro semanas finales de marcha blanca obligatoria ni sus criterios de aceptación \parencite[art.~17.3; pp.~12--13]{bases-admin}.
